\chapter*{Introduction}

In these notes, we will explore the fascinating world of Quantum Field Theory (QFT). QFT is a fundamental framework in physics that combines classical field theory, special relativity, and quantum mechanics. It provides a powerful language for describing the behavior of particles and fields at the most fundamental level. We assume are going to assume that the reader has a basic understanding of quantum mechanics and special relativity, as well as some familiarity with classical field theory and quantum field theory of free fields (Which will aniway be reviewed in the first chapter).

There exist different motivations for studying QFT: it provides a framework for describing the interactions between particles and fields, and it has been successful in explaining a wide range of phenomena, from the behavior of subatomic particles to the properties of condensed matter systems. Additionally:
\begin{itemize}
    \item QFT is the language of the Standard Model of particle physics, which describes the fundamental particles and their interactions (except for gravity); it is used for the quantization of systems with an \textbf{infinite number of degrees of freedom}, such as fields; it is essential for understanding the behavior of particles at high energies and small distances, where quantum effects become significant; and it provides a framework for describing the behavior of particles in curved spacetime, which is important for understanding phenomena such as black holes and the early universe.
    \item It is a successful framework for describing the behavior of particles and fields at the most fundamental level, incorporating both quantum mechanics and special relativity:
          \begin{itemize}
              \item At high energies we can witness \textbf{particles creation and annihilation}, which cannot be described by non-relativistic quantum mechanics; QFT provides a framework for describing these processes, in an infinite-dimensional system, where the number of particles is not fixed.
              \item QFT is a relativistic theory, which means that it is consistent with the principles of special relativity, implmenting \textbf{causality} and \textbf{Lorentz invariance}; this is essential for describing the behavior of particles at high energies and small distances, where relativistic effects become significant. So two events that are spacelike separated cannot influence each other
                    \[
                        (x-y)^2 = (x^0-y^0) - \sum_{i=1}^3 (x^i-y^i)^2 < 0 \implies [\mathcal{O}(x), \mathcal{O}(y)] = 0,
                    \]
                    thus the events are causally disconnected, ensuring causality.
          \end{itemize}
    \item Particles are described as excitations of underlying fields, which can be created and annihilated; this allows for a more complete description of particle interactions and the behavior of particles at high energies. Examples of particles and underlying fields include:
          \begin{itemize}
              \item \textbf{Scalar field}: \(\phi(x)\) \(\implies\) spin-0 particles (e.g., Higgs boson).
              \item \textbf{Spinor field}: \(\psi(x)\) \(\implies\) spin-1/2 particles (e.g., electrons, quarks).
              \item \textbf{Vector field}: \(A_\mu(x)\) \(\implies\) spin-1 particles (e.g., photons, W and Z bosons).
          \end{itemize}
    \item The \textbf{Standard Model} of particle physics is a quantum field theory that describes the fundamental particles and their interactions (except for gravity); it has been successful in explaining a wide range of phenomena, from the behavior of subatomic particles to the properties of condensed matter systems. It is a quantum field theory that describes the electromagnetic, weak, and strong interactions, through the exchange of gauge bosons among fields of spin 0, 1/2, and 1 particles.
    \item The inclusion of gravity in a quantum field theory framework is an open problem in theoretical physics, but can be approached by inserting the metric tensor of general relativity
          \[
              g^{\mu \nu} = \eta^{\mu \nu} + h^{\mu \nu}.
          \]
          While there are approaches to quantizing gravity, such as string theory and loop quantum gravity, a complete and consistent theory of quantum gravity has not yet been developed. Quanta of the gravitational field are called gravitons, which are hypothetical spin-2 particles that mediate the gravitational interaction, but they have not been observed experimentally. These notes will not cover quantum gravity, but it is an important area of research in theoretical physics which sometimes will be mentioned throughout the notes.
\end{itemize}

The main goal of these notes is to provide a comprehensive understanding of interacting quantum field theories, which are essential for describing the behavior of particles and fields at the most fundamental level, as well as mathematical methods for making theoretical predictions on interactions among particles. We will cover a wide range of topics, including:
\begin{itemize}\TODO{Check and change in the end}
    \item The quantization of free fields, which provides the foundation for understanding the behavior of particles and fields in quantum field theory.
    \item The construction of interacting quantum field theories, which are essential for describing the behavior of particles and fields at the most fundamental level.
    \item The use of perturbation theory and Feynman diagrams to calculate scattering amplitudes and other physical observables in quantum field theory.
    \item The renormalization of quantum field theories, which is necessary to make sense of the infinities that arise in perturbation theory and to extract meaningful physical predictions from the theory.
    \item The application of quantum field theory to the Standard Model of particle physics, which describes the fundamental particles and their interactions (except for gravity).
\end{itemize}

QFT is one of the most successful and powerful frameworks in physics, and it has been instrumental in our understanding of the fundamental nature of the universe. This theory has been tested and confirmed by a wide range of experiments, some of which tested the predictions of QFT to an extraordinary degree of precision. For example, the anomalous magnetic moment of the electron has been measured to an accuracy of one part in a trillion \((\sim 10^{-13})\), and the predictions of QFT have been confirmed to this level of precision. Additionally, the discovery of the Higgs boson at the Large Hadron Collider in 2012 provided further confirmation of the predictions of QFT and the Standard Model.

Another essential aspect of these notes is to provide a comprehensive understanding of the importance of symmetries in quantum field theory, which play a crucial role in determining the behavior of particles and fields. We will explore the concept of gauge symmetries, which are local symmetries that arise in quantum field theories and are responsible for the interactions between particles and fields. We will also discuss the role of global symmetries, which are symmetries that do not depend on spacetime coordinates, and their implications for the behavior of particles and fields. In the end we will argue that \textit{a QFT is essentially defined by its fields and its symmetries} (statement that will be made more precise in the end of the notes), and we will see how symmetries can be used to constrain the behavior of particles and fields, and to make predictions about their interactions.

For a more comprehensive introduction to quantum field theory, we recommend the following textbooks:\
\begin{itemize}
    \item \textbf{An Introduction to Quantum Field Theory}\TODO{check and change in the end} by Michael E. Peskin and Daniel V. Schroeder: This is a widely used textbook that provides a comprehensive introduction to quantum field theory, covering both the theoretical foundations and practical applications of the subject.
    \item \textbf{Quantum Field Theory} by Mark Srednicki: This textbook offers a clear and concise introduction to quantum field theory, with an emphasis on the physical intuition behind the mathematical formalism.
    \item \textbf{Quantum Field Theory in a Nutshell} by A. Zee: This book provides a more informal and intuitive introduction to quantum field theory, making it accessible to a wider audience while still covering essential topics.
    \item \textbf{The Quantum Theory of Fields} by Steven Weinberg: This three-volume series is a comprehensive and rigorous treatment of quantum field theory, written by one of the pioneers in the field. It is suitable for advanced students and researchers who want a deep understanding of the subject.
\end{itemize}

For more information about the notation and conventions used in these notes, please refer to the appendix, where we provide a detailed explanation of the symbols and mathematical conventions that we will use throughout the notes (Appendix \ref{app:notation}).

\vspace{1cm}
\textbf{Prerequisites:} \textit{Familiarity with the Lagrangian and Hamiltonian formalisms of Classical Mechanics, standard Quantum Mechanics, Special Relativity and the basics of Quantum Field Theory (free fields) is required.}